\subsection{EG gauge symmetry: BRST and Slavnov--Taylor consistency (strong-closure layer)}
\label{app:eg_brst_st_consistency}

\AuditTag This appendix upgrades the EFT BRST/Ward layer from Appendix~\ref{app:ward_brst_consistency} to a theorem-facing statement in the EG carrier (Appendix~\ref{app:eg_causal_perturbation_framework}).
The intent is to record a precise strong-closure target: \emph{existence (to finite order) of renormalized time-ordered products consistent with BRST/Slavnov--Taylor constraints}, with anomalies isolated as explicit cohomological obstructions.

\subsubsection{Classical BRST differential and interaction}
\paragraph{BRST differential.}
Let $s$ denote the classical BRST differential acting on the local functional algebra $\mathcal{F}_{\mathrm{loc}}$ (Appendix~\ref{app:ward_brst_consistency}).
We assume $s^2=0$ on $\mathcal{F}_{\mathrm{loc}}$ and that the classical interaction functional $V\in\mathcal{F}_{\mathrm{loc}}$ is BRST-invariant modulo total derivatives:
\[
sV=\partial_\mu J^\mu
\quad\text{(locally).}
\]

\paragraph{EFT input.}
\InterfaceTag The interaction $V$ is chosen from the finite EFT candidate family defined in Appendix~\ref{app:eft_representative_construction}, including gauge fields, ghosts, and the CAP-minimal scalar sector.

\subsubsection{Quantum BRST and Slavnov--Taylor operator}
\paragraph{Quantum BRST condition (target form).}
In the EG setting, BRST/Slavnov--Taylor constraints are encoded as normalization conditions on the time-ordered products $T_n$ (equivalently, on the S-matrix functional $S(V)$).
We record the strong-closure target in a form suitable for order-by-order theoremization:

\begin{theorem}[ST/BRST normalization in the causal carrier (finite order)]
\label{thm:eg_st_normalization_finite_order}
\MathTag Fix an EFT truncation and a target perturbative order.
Assume the anomaly filters in Appendix~\ref{app:anomaly_theorem_filters} hold for the declared gauge group and chiral content.
Then there exists a choice of EG normalization (extensions on diagonals) such that, up to the target order, the renormalized S-matrix $S(V)$ satisfies the Slavnov--Taylor identity.
\end{theorem}
\begin{proof}
This is a standard consequence of causal perturbation theory with gauge symmetry: locality/causal factorization reduces renormalization to finite local freedom at each order, and anomaly cancellation implies that the ST breaking term is cohomologically removable by local counterterms.
For references and an explicit order-by-order restoration formulation, see \cite{EpsteinGlaser1973,BrunettiFredenhagen2000Microlocal,Dutsch2019pAQFT,BecchiRouetStora1976,Tyutin1975} and Appendix~\ref{app:eg_st_restoration_algorithm}.
\end{proof}

\subsubsection{Cohomological obstruction: anomalies}
\paragraph{Anomaly as ST breaking.}
If ST cannot be imposed, the failure appears as a local breaking term $\Delta$ in the ST identity:
\[
\mathrm{ST}(S(V))=\Delta,
\]
where $\Delta$ is local, has controlled dimension/ghost number, and satisfies a consistency condition (Wess--Zumino type) induced by $s^2=0$.

\paragraph{Filter input.}
Appendix~\ref{app:anomaly_theorem_filters} provides the anomaly cancellation constraints for the Standard Model chiral content, including the global $SU(2)$ obstruction.
In the EG strong-closure chain, those constraints are treated as necessary conditions for $\Delta$ to be removable by finite renormalizations.

\subsubsection{Strong-closure theorem statement (finite order)}
\begin{theorem}[BRST/Slavnov--Taylor realizability to finite order (template)]
\label{thm:eg_st_realizability_template}
Fix an EFT truncation order $N$ and a finite-order expansion level in the EG construction.
Assume:
\begin{itemize}
\item the classical BRST differential $s$ is nilpotent on $\mathcal{F}_{\mathrm{loc}}$ and $V$ is BRST-invariant modulo total derivatives;
\item the local anomaly coefficients required by Appendix~\ref{app:anomaly_theorem_filters} vanish (and the global $SU(2)$ parity condition holds);
\item the EG scaling-degree/power-counting constraints are satisfied at the chosen truncation order.
\end{itemize}
Then there exists a choice of EG renormalization freedoms (extensions on the diagonals) such that the renormalized time-ordered products $T_n$ (equivalently $S(V)$) satisfy the ST identity up to the fixed order.
\end{theorem}

\AuditTag Theorem~\ref{thm:eg_st_realizability_template} is a theorem-facing statement whose carrier proof is provided by the causal/BRST machinery once the admissible local spaces are fixed.
In this manuscript, the admissible local counterterm space is normalized by the finite basis registry (Appendix~\ref{app:eg_normalization_freedom_basis}), and the inductive restoration is recorded algorithmically in Appendix~\ref{app:eg_st_restoration_algorithm}.

\subsubsection{Failure points (strong-closure)}
\AuditTag The strong-closure layer has explicit failure modes:
\begin{itemize}
\item If anomaly filters fail, ST cannot be restored by local counterterms (obstruction in BRST cohomology).
\item If power counting fails at the declared truncation, the finite-dimensional renormalization freedom claim breaks and the EFT truncation must be revised.
\item If the chosen observable class requires nonlocal composite operators beyond the EG control class, additional hypotheses are required (domain/regularity upgrade).
\end{itemize}

